\documentclass[10pt,a4paper]{amsart}
\usepackage[margin=19mm]{geometry}
\usepackage{amsmath,amssymb,mathtools}
\usepackage[scr=rsfs,cal=euler]{mathalfa}
\usepackage{xfrac}
\usepackage[unicode,hidelinks,bookmarks=false]{hyperref}
\newcommand{\ie}{i.e.\ }
\newcommand{\cf}{cf.\ }
\newcommand{\resp}{respectively\ }
\newcommand{\Subsection}{\subsection}
\providecommand{\nobreakdash}{\nobreak\hbox{-}}
\setlength{\emergencystretch}{2em}
\multlinegap0pt
\usepackage{mathtools}
\usepackage{tikz}
\newtheorem{proposition}{Proposition}
\newcommand{\N}{\mathbb{N}}
\renewcommand{\P}{\mathbb{P}}
\newcommand{\Reff}{\mathcal{R}_{\mathrm{eff}}}
\newcommand{\dG}{\operatorname{d}_{G}}
\title{Recurrence of strong-decay inhomogeneous long-range percolation clusters}
\author{Johannes B\"aumler, Lukas L\"uchtrath, Christian M\"onch}
\date{Source: arXiv:2608.25201v1 (CC BY 4.0)}
\begin{document}
\maketitle

\noindent\emph{Source and licence.} Recurrence of strong-decay inhomogeneous long-range percolation clusters. Version arXiv:2608.25201v1 (CC BY 4.0), distributed by arXiv under the Creative Commons Attribution 4.0 International licence, http://creativecommons.org/licenses/by/4.0/, as recorded in the arXiv licence metadata for this exact version and shown on the abstract page https://arxiv.org/abs/2608.25201v1. Full licence notice: \texttt{licenses/arxiv-2608.25201-CC-BY-4.0.txt}.\par\smallskip

\noindent\textit{Editorial scope.} Counted unit: the article's proposition prop:fixed-box-distance-rec (source0.tex lines 597--617). The article's complete original proof of it is delivered with it (source0.tex lines 619--727, one \texttt{proof} environment of 109 lines). No mathematics is written, repaired or re-derived: the statement and the proof are the article's own lines, in the article's own notation, and no retained line is substituted or re-flowed.  No further source-local context is needed: every cross-reference of every retained line resolves inside the two delivered spans.  Nothing was omitted from any retained span, so the package carries no omission record.  No retained line contains a citation, so the wrapper carries no bibliography and no bibliography entry.  No retained line contains a figure environment, an image-inclusion command or drawing code, so no image is delivered and the package declares no asset dependency.  Editorial formatting only, in addition: an AMS \texttt{amsart} preamble that restates the article's own declarations and macros used by the retained lines, namely the environments \texttt{proposition} on separate counters, together with the standard packages those lines need.  The retained environments print in this document as Proposition 1 in order of appearance, whereas the article prints the same items with the numbering of its own full document.  The complete native source of the article as distributed in the arXiv e-print for this exact version is bundled unchanged as \texttt{sources/208555/source0.tex}.
\begin{proposition}
\label{prop:fixed-box-distance-rec}
Suppose that \(G\) is stationary, ergodic, and locally finite, that its degree
measure has finite intensity, and that for all $L \geq 0$ there exists \(\eta>0\) such that
\begin{equation}
\label{eq:fixed-box-distance}
  \P\big(\mathcal D(R)\big)\longrightarrow1
  \qquad\text{as }R\to\infty,
\end{equation}
where
\[
  \mathcal D(R)\coloneqq 
  \left\{
    \dG(v,w)\ge\eta|v-w|
    \text{ for all }v\in V\cap Q_L(0),\
    w\in V\cap Q_R(0)^{\mathsf c}
  \right\}.
\]
Then almost surely every connected component of
\(G\) is recurrent.
\end{proposition}

\begin{proof}
If \(\theta_M=0\), then the non-negative random variable
\(T(V\cap[0,1)^2)\) vanishes almost surely. By stationarity, the same is true
in every integer translate of the unit box, so every vertex has degree zero
almost surely and the claim is immediate. We may therefore assume that
\(\theta_M>0\).

Fix \(L\in\N\), and let \(\eta>0\) be as in
\eqref{eq:fixed-box-distance}. We prove the claim for all components meeting
\(Q_L(0)\); a countable intersection over \(L\) will then give the result.
Set \(R_n=4^n\).  Since
\(\P(\mathcal D(R_n))\to1\), we get that
\[
  \P\big(\mathcal D(R_n)\text{ infinitely often}\big) = \P\left( \bigcap_{m=1}^{\infty} \bigcup_{n= m}^{\infty} \mathcal D(R_n)\right)
  \geq 
  \lim_{m\to \infty} \sup_{n \geq m}
  \P\left( \mathcal D(R_n)\right)=1.
\]

Let \(M(A)=T(V\cap A)\) be the degree measure.  The multiparameter ergodic
theorem gives
\[
  \frac{M(Q_R(0))}{R^2}\longrightarrow\theta_M
  \qquad\text{almost surely, for } R\to \infty.
\]
Thus we see that there exist, almost surely, infinitely many $n \in \N$ so that
the event \(\mathcal D(R_n)\) occurs and
\begin{equation}
\label{eq:direct-degree-volume}
  M(Q_{R_n}(0))\le 2 \theta_M R_n^2.
\end{equation}


Work on this event and fix \(v\in V\cap Q_L(0)\). If \(\mathcal D(R_n)\)
occurs and \(R_n\ge2L\), then every
\(w\notin Q_{R_n}(0)\) satisfies
\[
  |v-w|
  \ge \frac{R_n-L}{2}
  \ge \frac{R_n}{4},
\]
and therefore
\begin{equation}
\label{eq:direct-distance-outside}
  \dG(v,w)\ge\frac{\eta R_n}{4}.
\end{equation}

Put \(a=\eta/16\).  For \(k\ge0\), define
\[
  D_k(v) \coloneqq \{w\in V:\dG(v,w)\le k\},
\]
and, for \(k\ge1\), let
\[
  F_k(v) \coloneqq
  \big\{\{x,y\}\in E:|\{x,y\}\cap D_{k-1}(v)|=1\big\}.
\]
Local finiteness implies that every \(D_k(v)\), and hence every \(F_k(v)\),
is finite.  If the component of \(v\) is infinite, each \(F_k(v)\) is a cutset
separating \(v\) from infinity.  The cutsets are pairwise edge-disjoint: the
graph distances from \(v\) of the endpoints of an edge differ by at most one.

For a scale at which \(\mathcal D(R_n)\) occurs, consider
\[
  I_n\coloneqq 
  \big\{k\in\N:
    \lceil aR_n\rceil\le k\le\lfloor2aR_n\rfloor
  \big\}.
\]
By \eqref{eq:direct-distance-outside},
\(D_k(v)\subseteq Q_{R_n}(0)\) for every \(k\in I_n\).  Thus, all edges in
\(\bigcup_{k\in I_n}F_k(v)\) have both endpoints in \(Q_{R_n}(0)\).  Since
these edge sets are disjoint, \eqref{eq:direct-degree-volume} yields
\[
  \sum_{k\in I_n}|F_k(v)|
  \le M(Q_{R_n}(0))
  \le 2 \theta_M R_n^2.
\]
Moreover, \(|I_n|\ge aR_n/2\) for all large enough \(n\). Let $N$ be large enough so that $ M(Q_{R_n}(0))\le 2 \theta_M R_n^2$ and \(|I_n|\ge aR_n/2\) for all $n\geq N$.  If some \(F_k(v)\) is
empty, then \(D_{k-1}(v)\) is a finite set with no edge leaving it, so the
component of \(v\) is finite and hence recurrent.  Otherwise,
Cauchy--Schwarz gives
\begin{equation}
\label{eq:direct-resistance-band}
  \sum_{k\in I_n}\frac1{|F_k(v)|}
  \ge
  \frac{|I_n|^2}{\sum_{k\in I_n}|F_k(v)|}
  \ge
  \frac{a^2}{8 \theta_M}
\end{equation}
for $n\geq N$.
Because \(R_{n+1}=4R_n\), the bands \(I_n\), \(n \in \N\), are pairwise disjoint for all
large \(n\).  Infinitely many of the corresponding distance events occur, so
collecting their cutsets and applying Nash--Williams gives
\[
  \Reff(v,\infty)
  \ge
  \sum_{\substack{n \geq N: \\ \mathcal D(R_n)\text{ occurs}}}
  \ \sum_{k\in I_n}\frac1{|F_k(v)|}
  \ge
  \sum_{\substack{n \geq N: \\ \mathcal D(R_n)\text{ occurs}}}
  \frac{a^2}{8 \theta_M}
  =
  \infty.
\]
Thus the component of \(v\) is recurrent. Since $v$ was arbitrary in
\(V\cap Q_L(0)\), every component meeting this box is recurrent on the same
probability-one event. Intersecting these events over \(L\in\N\) proves the
claim for all components.
\end{proof}

\end{document}
